\section{Experiment 3: Simulation of UNIX Commands}
\label{exp:simulate}

\subsection*{Aim}
To write C programs that simulate the behaviour of the UNIX commands \texttt{ls} and \texttt{grep}.

\subsection*{Theory}
UNIX commands are ordinary programs that use system calls and library functions. The \texttt{ls} command reads a directory and prints its entries; a user program can reproduce this by calling \texttt{opendir} and \texttt{readdir}. The \texttt{grep} command reads a file line by line and prints the lines that contain a given pattern; a user program can reproduce this with \texttt{fgets} and \texttt{strstr}.

The simulations in this experiment are deliberately small. They reproduce the core visible behaviour of the commands, not every option of the real utilities.

\subsection{Simulation of the \texttt{ls} Command}

\paragraph{Aim.}
To list the contents of a directory supplied as a command-line argument.

\paragraph{Algorithm.}
\begin{enumerate}
    \item Start the program and check that exactly one argument (the directory path) is supplied.
    \item Open the directory with \texttt{opendir(argv[1])}. If it fails, print an error and exit.
    \item Print a heading containing the directory name.
    \item Call \texttt{readdir} repeatedly and print the name of each entry.
    \item Close the directory stream with \texttt{closedir}.
    \item Stop the program.
\end{enumerate}

\paragraph{Program.}
\lstinputlisting[style=cstyle, caption={Simulation of the \texttt{ls} command (\texttt{expt\_3\_1.c})}]{Programs/expt_3_1.c}

\paragraph{Sample Output.}
\begin{lstlisting}[style=outputstyle, caption={Listing a directory with the simulated \texttt{ls}}]
$ gcc expt_3_1.c -o expt_3_1
$ ./expt_3_1 .
Contents of .:
.
..
01_fibonacci.sh
02_palindrome.sh
03_power.sh
04_second_largest.sh
05_prime.sh
06_sum_of_digits.sh
07_multiplication_table.sh
08_even_or_odd.sh
09_leap_year.sh
10_factorial.sh
11_swap_two_integers.sh
expt_3_1.c
expt_3_1
\end{lstlisting}

\subsection{Simulation of the \texttt{grep} Command}

\paragraph{Aim.}
To search for a word in a file and print every matching line together with its line number.

\paragraph{Algorithm.}
\begin{enumerate}
    \item Start the program and check that a file name and a search word are supplied.
    \item Open the file in read mode. If it cannot be opened, print an error and exit.
    \item Read each line with \texttt{fgets} and increment the line counter.
    \item Remove the trailing newline from the line.
    \item Use \texttt{strstr(line, word)} to test whether the line contains the word.
    \item If a match is found, print the file name, line number, and the line.
    \item Close the file and stop the program.
\end{enumerate}

\paragraph{Program.}
\lstinputlisting[style=cstyle, caption={Simulation of the \texttt{grep} command (\texttt{expt\_3\_2.c})}]{Programs/expt_3_2.c}

\paragraph{Sample Output.}
\begin{lstlisting}[style=outputstyle, caption={Searching for the word \texttt{Fibonacci} in a file}]
$ gcc expt_3_2.c -o expt_3_2
$ ./expt_3_2 01_fibonacci.sh Fibonacci
01_fibonacci.sh:3: # Write a script to print Fibonacci series up to n
terms
01_fibonacci.sh:11: echo "Fibonacci Series:"
\end{lstlisting}

\subsection{Counting Occurrences with a \texttt{grep}-style Program}

\paragraph{Aim.}
To search for a pattern in a file, print every matching line, and count the total number of occurrences of the pattern.

\paragraph{Algorithm.}
\begin{enumerate}
    \item Start the program and read the file name and the pattern from the command line.
    \item Open the file and initialise the line number and the total occurrence count to zero.
    \item Read each line with \texttt{fgets} and increment the line number.
    \item Repeatedly apply \texttt{strstr} to the line, advancing the search pointer after each match, and count the non-overlapping occurrences.
    \item If the current line contains at least one match, print the line with its number.
    \item Add the line's occurrences to the total and continue until end of file.
    \item Print the total number of occurrences, close the file, and stop.
\end{enumerate}

\paragraph{Program.}
\lstinputlisting[style=cstyle, caption={Occurrence counting with a \texttt{grep}-style program (\texttt{expt\_3\_3.c})}]{Programs/expt_3_3.c}

\paragraph{Sample Output.}
\begin{lstlisting}[style=outputstyle, caption={Counting occurrences of the pattern \texttt{is}}]
$ gcc expt_3_3.c -o expt_3_3
$ ./expt_3_3 02_palindrome.sh is
02_palindrome.sh:3: # Write a script to check if a string is a
palindrome
02_palindrome.sh:12:     echo "The string is a palindrome."
02_palindrome.sh:14:     echo "The string is not a palindrome."
Total occurrences of "is": 3
\end{lstlisting}

\subsection*{Result}
The UNIX commands \texttt{ls} and \texttt{grep} were simulated in C. The programs listed a directory, printed matching lines with line numbers, and counted the total number of pattern occurrences successfully.
